% BEGIN GENERATED ARTIFACT: prf:tautological-implication-valid-argument-form-propositional-logic
\newpage
\phantomsection
\label{prf:tautological-implication-valid-argument-form-propositional-logic}
\LRAProofFor{thm:tautological-implication-valid-argument-form-propositional-logic}

\begin{remark*}[Return]
\hyperref[thm:tautological-implication-valid-argument-form-propositional-logic]{Return to Theorem (Tautological Implication Gives a Valid Argument Form).}
\end{remark*}

\begin{theorem*}[Tautological Implication Gives a Valid Argument Form]
Let \(\Gamma=\{\gamma_1,\ldots,\gamma_n\}\subseteq\WFF\) and let
\(\varphi\in\WFF\). If
\[
(\gamma_1\land\cdots\land\gamma_n)\to\varphi
\]
is a tautology, then the argument form \((\Gamma,\varphi)\) is valid.
\end{theorem*}

\begin{proof}
\textbf{Professional Standard Proof.}~
TODO: supply the compact proof.
\end{proof}

\begin{proof}
\textbf{Detailed Learning Proof.}~
TODO: supply the detailed learning proof.
\end{proof}

\begin{remark*}[Proof structure]
\end{remark*}

\begin{dependencies}
\begin{itemize}
  \item \hyperref[def:tautology-propositional-logic]{Tautology}
  \item \hyperref[def:valid-argument-form-propositional-logic]{Valid Argument Form}
  \item \hyperref[def:logical-consequence-propositional-logic]{Logical Consequence}
\end{itemize}
\end{dependencies}

\clearpage
% END GENERATED ARTIFACT: prf:tautological-implication-valid-argument-form-propositional-logic
